\documentclass[11pt]{article}

\usepackage[letterpaper,margin=1in]{geometry}
\usepackage{microtype}
\usepackage{amsmath}
\usepackage{amssymb}
\usepackage{booktabs}
\usepackage{array}
\usepackage[table,dvipsnames]{xcolor}
\usepackage{enumitem}
\usepackage[most]{tcolorbox}
\usepackage{float}
\usepackage[hidelinks]{hyperref}

\emergencystretch=3em

% Plain-language callout.
\newtcolorbox{plain}{
  enhanced, breakable,
  colback=Sepia!7, colframe=Sepia!55, boxrule=0.4pt, arc=2pt,
  left=7pt, right=7pt, top=4pt, bottom=4pt,
  fontupper=\small,
  coltitle=Sepia!72!black, fonttitle=\bfseries\sffamily\footnotesize,
  title={IN PLAIN TERMS}
}
% Definition callout.
\newtcolorbox{definitionbox}{
  enhanced, breakable,
  colback=infoblue!4, colframe=infoblue!45, boxrule=0.4pt, arc=2pt,
  left=7pt, right=7pt, top=4pt, bottom=4pt,
  fontupper=\small,
  coltitle=infoblue!80!black, fonttitle=\bfseries\sffamily\footnotesize,
  title={DEFINITION}
}

\definecolor{crit}{HTML}{C0392B}
\definecolor{high}{HTML}{E67E22}
\definecolor{med}{HTML}{F1C40F}
\definecolor{low}{HTML}{27AE60}
\definecolor{infoblue}{HTML}{2C3E50}
\definecolor{rule}{HTML}{BDC3C7}

\newcommand{\code}[1]{\texttt{#1}}
\newcommand{\rfckw}[1]{\textsc{#1}}

\title{\bfseries PAIN Relief: A Verified-Control Method for\\
Reducing FedRAMP Rev5 VDR/VER PAIN and Exploitability}
\author{Matthew Venne \\ \small Chief Technology Officer, stackArmor}
\date{July 2026}

\begin{document}
\maketitle

\begin{abstract}
\noindent
The companion memo \emph{A Deterministic, CVSS--Environmental Method for FedRAMP
Rev5 VDR/VER Vulnerability Prioritization} derives a finding's Potential Agency
Impact (PAIN) and its remediation deadline. That memo also establishes a
mitigation \emph{ladder}: PAIN comes down when cited, signed evidence lowers a
Modified impact metric. This paper makes the ladder systematic. It defines
\textbf{PAIN Relief}: a governed, versioned catalog that pairs a
\emph{machine-verified} security control with the specific weakness class
(by CWE) it provably counters and the deterministic scoring reduction that
pairing earns. Two levers move: verified controls that bound what a successful
exploit can do lower a Modified impact metric and step the PAIN level down;
verified controls that reduce the probability of successful exploitation lower a
\emph{local} estimate of that probability and can move a finding from
likely-exploitable to not. Neither lever mutates a published input --- not the
CVSS base vector, not the EPSS score, not the CISA KEV clock. The method's
guiding line is \emph{enforced-and-machine-verifiable}, not \emph{claimed}: a
control earns relief only when enforcement infrastructure applies it and a
collector can read it. The catalog is offered as an example; each provider
governs its own.
\end{abstract}

\tableofcontents

\section{Status and Relationship to the Companion Papers}
This document is informational. It does not define a FedRAMP requirement; it
specifies a method for lowering the PAIN and exploitability of a finding within
the existing VDR/VER rules, building directly on the PAIN memo's mitigation
ladder. It consumes the internet-reachability determination of \emph{An
Exclusion-Gated Method for Determining Internet Reachability} and the disposition
vocabulary of \emph{A CycloneDX VEX Profile for FedRAMP Rev5 VDR/VER}. The key
words \rfckw{must}, \rfckw{should}, and \rfckw{may} follow RFC~2119.

\section{The Problem: Good Hygiene Must Pay}
A provider that runs its workloads with least privilege, deny-by-default egress,
read-only filesystems, and enforced runtime protection is measurably harder to
harm than one that does not. Yet a raw scan scores an identical CVE identically
on both. If secure configuration changes no number a provider must report, it is
all cost and no credit, and the rational provider treats hardening as a burden.
The PAIN method already refuses that outcome for one input --- asset criticality,
through the CVSS Environmental requirements --- by re-weighting a vulnerability's
impact by what is at stake on the specific asset. PAIN Relief extends the same
principle from \emph{what the asset is} to \emph{what protects it}.

\begin{plain}
A scanner scores a flaw as if every system running it were equally exposed. Two
providers run the same flaw; one has locked the system down and one has not.
PAIN Relief is how the locked-down provider's work shows up in the score ---
lower impact, or a longer clock --- so that doing security well is worth
something on paper, not just in principle.
\end{plain}

\subsection{The axis that matters: enforced, not claimed}
The tempting mistake is to reward controls a provider \emph{asserts}. Input
sanitization, a WAF signature set, ``we validate everything'' --- these are
claims about behavior, unverifiable at assertion time and silently perishable on
the next refactor. PAIN Relief admits a control only when it is
\textbf{machine-verifiable} --- read from enforced configuration or telemetry a
collector can inspect (a NetworkPolicy, a \code{securityContext}, an RDS flag) ---
or carried by a signed, dated, reusable attestation for structure a platform
cannot expose. A control whose effectiveness is a claim stays in the VEX
disposition lane of the companion profile, decided per finding by a person. The
line is not \emph{a priori} versus \emph{a posteriori}; it is
\emph{enforced-and-verifiable} versus \emph{asserted}.

\section{Two Levers}
\label{sec:levers}
PAIN Relief moves exactly two things, and keeps them strictly apart. A control
belongs to one lever, chosen by its causal mechanism: a control that bounds
\emph{what a successful exploit yields} moves impact; a control that reduces
\emph{whether the exploit succeeds} moves exploitability. Because the assignment
follows the mechanism, a control never scores on both axes, and nothing
double-counts.

\subsection{Impact: lowering a Modified metric}
\label{sec:impact}
The PAIN memo scores a finding through the CVSS Modified Impact Sub-Score, built
from the per-dimension impacts $C,I,A$ and the asset's requirements. A verified
control that provably counters a weakness class lowers the affected Modified
impact metric from High to Low before that score is computed:
\begin{equation}
\text{High} \;\longrightarrow\; \text{Low}, \qquad \text{never } \text{None}.
\label{eq:rung}
\end{equation}
The floor is deliberate. A claim of \emph{zero} impact is an applicability claim,
and it travels as a VEX disposition (\code{code\_not\_present},
\code{code\_not\_reachable}), not as relief. The reduced $C$, $I$, or $A$ flows
through the memo's existing arithmetic unchanged; the PAIN level and deadline
follow. There is one rung, and it does not stack: several controls that each
argue a dimension down still lower it by a single step, with every contributing
control named in the finding's evidence.

\subsection{Exploitability: lowering a local estimate}
\label{sec:exploit}
FedRAMP's exploitability axis is binary: a finding is likely-exploitable (LEV) or
not (NLEV), and the two select different remediation columns. The memo computes
\[
\text{LEV} \;=\; \text{KEV} \ \lor\ (\text{EPSS} \ge \tau)\ \lor\ \text{floor},
\qquad \tau = 0.70,
\]
where the floor is the \code{FRD-LEV} rule that an automated unauthenticated
internet exploit is likely-exploitable. PAIN Relief never edits the published
EPSS score --- a global prediction is not a provider's to rewrite --- just as the
impact lever never edits the CVSS base vector. Instead a verified control lowers
a \emph{local} estimate of exploitation probability, and the binary threshold
does the rest:
\begin{equation}
\mathrm{adjustedEPSS} \;=\; \max\!\Big(\ \mathrm{EPSS}\times\!\!\prod_{r\in R}\rho_r,\ \
\mathrm{EPSS}\times\phi\ \Big),
\label{eq:adjepss}
\end{equation}
where $R$ is the set of applicable verified likelihood controls, each carrying a
\emph{residual factor} $\rho_r\in(0,1)$ --- the share of exploitation risk it
leaves behind --- and $\phi$ is a governed floor on total reduction. LEV is then
recomputed with \code{adjustedEPSS} in place of EPSS.

\begin{definitionbox}
$\rho_r$ (\emph{residual factor}): a control that halves exploitation likelihood
carries $\rho=0.5$; one that trims it modestly carries $\rho=0.85$. Lower is
stronger. $\phi$ (\emph{stacking floor}): the most that stacking may reduce the
estimate, e.g.\ $\phi=0.5$ caps total reduction at one half. Both are governed,
back-tested constants, documented and challengeable like the PAIN cut points.
\end{definitionbox}

Three properties of Equation~\eqref{eq:adjepss} matter. First, \textbf{it stacks
multiplicatively}: independent controls compound, so layered runtime defense is
rewarded --- this is the one place defense-in-depth earns more than any single
control. Second, \textbf{the threshold self-calibrates}: a modest control drops a
marginal finding ($0.72\times0.85=0.61$, now NLEV) but not a near-certain one
($0.99\times0.85=0.84$, still LEV), which is exactly right --- one control should
not defuse a vulnerability being exploited everywhere. Third, and absolutely,
\textbf{KEV is frozen}: a Known Exploited Vulnerability stays LEV regardless of
$\rho$, and its BOD~26-04 clock is untouched. The \code{FRD-LEV} floor is a
reachability premise, not a probability, so it is defeated only in the binary ---
by a qualifying edge-authentication boundary --- never by a residual factor.

\begin{plain}
EPSS is a worldwide guess at how likely a flaw is to be attacked. We never change
that number. We keep a second, local number that starts at EPSS and drops when we
can verify a control that makes the attack less likely to work here. If the local
number falls below the line, the finding earns the slower clock. Stacked controls
push it lower together --- but a flaw that is already being exploited in the wild
(KEV) never moves, and the change is always shown beside the original so nothing
is hidden.
\end{plain}

\section{The Catalog}
\label{sec:catalog}
A PAIN Relief entry is a row pairing a control, the CWE class it counters, and
the move it earns. Rows are keyed to CWE because the finding carries a CWE and the
control's causal story is a statement about a weakness class --- ``deny-by-default
egress breaks the second-stage fetch of injection and staged-execution flaws''
is a claim about \code{CWE-917}, \code{CWE-94}, and their kin, not about a single
CVE. Where a control bounds what \emph{any} code execution yields regardless of
how it began, the row keys to an outcome class (all execution-capable weaknesses,
including memory corruption) rather than an enumerated list. A finding with no
CWE, or an unrecognized one, earns nothing: relief fails open toward the higher
score.

\begin{table}[H]
\centering
\renewcommand{\arraystretch}{1.25}
\footnotesize
\setlength{\tabcolsep}{5pt}
\begin{tabular}{l l l l}
\toprule
\textbf{Verified control} & \textbf{Counters (CWE)} & \textbf{Lever} & \textbf{Move} \\
\midrule
Egress default-deny        & 917, 94, 98, 829 & impact & MC/MI High$\to$Low \\
Read-only root filesystem  & 22, 434, 73       & impact & MI High$\to$Low \\
No shell in image          & 78, 77            & impact & MC/MI High$\to$Low \\
Least-privilege DB grant   & 89                & impact & MC/MI High$\to$Low \\
Verified HA + rate limit   & crash-class DoS   & impact & MA High$\to$Low \\
Blocking-mode EDR/RASP     & (any)             & likelihood & $\rho=0.85$ \\
Edge auth (remote-access)  & (any)             & likelihood & defeats \code{FRD-LEV} floor \\
\bottomrule
\end{tabular}
\caption{Illustrative PAIN Relief rows. The maintained catalog is proprietary and
provider-governed; this sample is an existence proof, not a standard --- exactly
the posture of the archetype catalog in the companion memo. Every applied credit
ships its row identifier, rationale, and the catalog version in the finding's
evidence.}
\label{tab:catalog}
\end{table}

\section{Governance}
\label{sec:governance}
Adding a row grants a discount across the estate, so a row is a security setting
and is reviewed like one. Each row \rfckw{must} satisfy:
\begin{itemize}[leftmargin=1.5em, nosep]
\item \textbf{A causal story, not an adjective.} ``Egress-deny breaks the staged
  fetch'' qualifies; ``EDR makes exploitation harder'' does not. The row names the
  mechanism by which the control interrupts the weakness class.
\item \textbf{Machine-verifiable or attested, never claimed.} The control is
  proven from enforced configuration, telemetry, or a signed reusable artifact.
  This is the clause that keeps input sanitization and WAF signature coverage
  \emph{out} of the catalog and in the VEX lane: both are bets that a pattern
  matches every variant an adversary can invent, and Log4Shell's same-week
  obfuscations are the standing proof they are not verifiable facts.
\item \textbf{CWE-keyed and fail-open.} A row applies only to the classes it
  lists; missing or generic CWEs earn nothing.
\item \textbf{One rung; never None.} Impact moves High to Low only; zero-impact is
  a VEX applicability claim. Exploitability moves the binary via
  Equation~\eqref{eq:adjepss}; nothing here touches reachability or KEV dates.
\item \textbf{Versioned, pinned, and evidenced.} A score is reproducible only
  against a named catalog release; the row id and version travel in every
  evidence line. Refusals are recorded, so the catalog is curated, not
  accumulated.
\item \textbf{Back-tested before adoption.} Residual factors, the stacking floor,
  and any new row are calibrated against real findings; a row that moves a large
  fraction of the estate is presumed miscalibrated until argued otherwise.
\end{itemize}

\section{What PAIN Relief Never Does}
No entry mutates the published EPSS score or the vendor CVSS base vector; both the
original and the adjusted values appear in output. No entry alters the
internet-reachability determination --- that is the companion exclusion-gated
method's job, and PAIN Relief merely consumes its result. No entry extends or
softens a CISA KEV / BOD~26-04 remediation date; a Known Exploited Vulnerability
carries its clock untouched no matter how thoroughly it is mitigated. And no
entry lowers a Modified impact dimension to None; that stronger claim is a signed
VEX disposition, separately governed.

\section{Worked Examples}
\label{sec:examples}

\paragraph{Impact --- egress-deny on a Log4Shell-class finding.} A
\code{CWE-917} finding scores $I=$ High on an application asset, placing it at,
say, N4. The workload runs under a verified default-deny egress NetworkPolicy
with no wildcard allow. Log4Shell's remote-code-execution requires an outbound
fetch of a second-stage payload; the egress policy removes the outbound leg, so
the injection still fires but cannot complete. The egress-deny row lowers Modified
Integrity High$\to$Low; recomputed through the memo's arithmetic the finding lands
at N3 --- a longer deadline, earned by a control the plugin read directly from the
cluster. The reachability answer is unchanged: the finding is still
internet-reachable; what changed is what it can accomplish.

\paragraph{Exploitability --- blocking EDR crossing the threshold.} A finding
carries EPSS $0.74$ (LEV, since $0.74 \ge 0.70$), is not KEV, and the floor does
not apply. The workload runs blocking-mode EDR in enforcing mode, verified from
its policy export, carrying $\rho = 0.85$. Then
$\mathrm{adjustedEPSS} = 0.74 \times 0.85 = 0.63 < 0.70$, and the finding moves to
NLEV --- the slower remediation column. Had EPSS been $0.99$, the same control
would leave it at $0.84$, still LEV: the threshold ensures a single control
relieves the marginal case, not the near-certain one. Published EPSS remains
$0.74$ in the record beside the adjusted value.

\begin{plain}
Same two flaws, two providers. On the hardened estate the first drops a PAIN
level and the second drops onto a slower clock, purely because a scanner-adjacent
tool could \emph{see} the controls that make those flaws less dangerous here. On
the unhardened estate, nothing moves. That difference is the entire point.
\end{plain}

\section{Architecture: Capture and Evaluation Are Separate}
The determination splits across a clean boundary. A collector with cluster
access --- the Kubernetes plugin --- \emph{captures} only: the finding's CWEs and
a neutral block of raw Kubernetes-observed facts per workload
(\code{securityContext} settings, NetworkPolicy rules, replica and disruption
posture, service-account and volume facts). It makes no interpretation: it emits
that a NetworkPolicy denies egress to a link-local address, not that ``IMDS is
protected.'' A downstream evaluator, which need not touch the cluster, loads the
governed catalog, maps the captured facts onto the control predicates, and
applies Equations~\eqref{eq:rung} and~\eqref{eq:adjepss}. The split keeps the
privileged collector simple and auditable and lets the catalog --- the part that
embodies judgment --- evolve under governed review without redeploying the
collector.

\section{Conclusion}
PAIN Relief closes the loop the PAIN memo opened. The memo re-weights a
vulnerability by what is at stake on the asset; this paper re-weights it by the
controls that verifiably protect the asset, using the same Environmental
machinery and the same discipline of citing evidence. Impact steps down a rung
when a verified control bounds the damage; exploitability crosses a threshold when
verified controls make success less likely; and neither ever edits a published
input, softens a KEV clock, or turns a claim into a discount. The result gives a
provider what raw scanning never has: a deterministic, auditable reason that good
security configuration lowers the number --- and therefore a reason to build
securely in the first place.
\end{document}
